\documentclass[../main.tex]{subfiles}
\graphicspath{{\subfix{../images/}}{\subfix{../tables_figures_sync/}}}

\begin{document}
\section{Abstract}
\textbf{Background:} As climate change accelerates, the frequency and severity of climate disasters in the United States has also increased. In order to respond effectively, we need a better understanding of both the effects of climate disasters themselves as well as the effects of our interventions -- in both directions. 

One recent intervention that started in California and has grown in popularity around the United States are Public Safety Power Shutoffs (PSPS). However, similar to other interventions, often this preventative measure ends up occurring simultaneously with wildfire \PM nonetheless. The objective of this study was to understand the effects of wildfire \PM alone, PSPS events alone, and the co-occurrence of wildfire \PM and PSPS events.

\textbf{Methods:} We conducted a \textit{time-stratified case-crossover} study among 1.2 million adults in California from 2017--2022. We used conditional logistic regression to estimate \texttt{odds ratios} (ORs) and 95\% confidence intervals (CIs).

\textbf{Results:} We observed a 15\% increase in respiratory emergency department visits (OR: 1.15, 95\% CI: 1.08--1.22) on ``high smoke'' days compared to matched control days.

\textbf{Conclusions:} Wildfire smoke exposure is associated with increased respiratory healthcare utilization, highlighting the need for \underline{public health interventions during smoke events}.

% ============================================================
% FORMATTING EXAMPLES EXPLAINED:
% ============================================================
% \textbf{text}     = bold text
% \textit{text}     = italic text  
% \underline{text}  = underlined text
% \texttt{text}     = code-style text
% \PM               = custom command for PM2.5 (defined in doc_setup.tex)
% --                = en-dash (for ranges like 2017--2022)
% ---               = em-dash (for breaks in thought---like this)
% ``text''          = proper quotation marks (not "text")
% \%                = percent sign (% alone starts a comment)
% ============================================================

\clearpage
\end{document}